\documentclass{ieeeaccess}
\usepackage{cite}
\usepackage{amsmath,amssymb,amsfonts}
\usepackage{algorithmic}
\usepackage{algorithm}
\usepackage{graphicx}
\usepackage{textcomp}
\def\BibTeX{{\rm B\kern-.05em{\sc i\kern-.025em b}\kern-.08em
    T\kern-.1667em\lower.7ex\hbox{E}\kern-.125emX}}
\begin{document}
\history{Date of publication xxxx 00, 0000, date of current version xxxx 00, 0000.}
\doi{10.1109/ACCESS.2017.DOI}

\title{VegShift: Detecting Crop Viability Loss Events After Climate Zone Transitions in Indian Cities Using Multi-Source Geospatial and Temporal Deep Learning}

\author{\uppercase{Author Name}\authorrefmark{1}}
\address[1]{Department, Institution, City, Country (e-mail: email@institution.ac.in)}

\tfootnote{Placeholder for funding and support acknowledgment.}

\markboth
{Author \headeretal: VegShift: Detecting Crop Viability Loss Events}
{Author \headeretal: VegShift: Detecting Crop Viability Loss Events}

\corresp{Corresponding author: Author Name (e-mail: email@institution.ac.in).}

\begin{abstract}
Placeholder abstract: This paper presents VegShift, a framework for detecting Crop Viability Loss Events (CVLEs) triggered by Köppen-Geiger climate zone transitions across ten major Indian cities over 2000–2024, using a multi-source fusion of atmospheric, hydrological, and agro-ecological data with an eight-model comparative deep learning study.
\end{abstract}

\begin{keywords}
crop viability loss, climate zone transition, Köppen-Geiger, groundwater depletion, temporal deep learning, TFT, CVLE, India agriculture
\end{keywords}

\titlepgskip=-15pt

\maketitle

% -----------------------------------------------------------------------
\section{Introduction}
\label{sec:introduction}
% -----------------------------------------------------------------------

\PARstart{C}{limate-driven} agricultural risk in India remains poorly quantified at the city scale, particularly the moment when a shifting climate permanently renders a crop non-viable. Existing work either uses static agroclimatic zones or single-variable thresholds; VegShift addresses this gap by formally timestamping viability loss events tied to climate zone transitions. This paper makes three contributions: (1) the CVLE dual-deficit event definition, (2) a 3-source geospatial dataset fusion spanning 25 years across 10 cities, and (3) a comparative study of 8 models with uncertainty quantification.

% -----------------------------------------------------------------------
\section{Related Work}
\label{sec:related_work}
% -----------------------------------------------------------------------

\subsection{Köppen-Geiger Classification in Agriculture}
Placeholder: Prior use of Köppen zones for agroclimatic zoning and its limitations under non-stationary climate.

\subsection{Crop Failure Prediction}
Placeholder: Survey of ML-based crop failure prediction methods, noting the predominance of single-variable rainfall or temperature thresholds.

\subsection{Groundwater Depletion and Agricultural Interaction}
Placeholder: Literature on CGWB-documented aquifer depletion trends and their link to irrigation-dependent crop failure in India.

\subsection{Temporal Deep Learning for Climate Applications}
Placeholder: Review of TFT, LSTM, TCN, and Transformer architectures applied to climate time-series classification and forecasting.

% -----------------------------------------------------------------------
\section{Study Area and Data}
\label{sec:data}
% -----------------------------------------------------------------------

\subsection{Study Cities and Crop Assignments}
Placeholder: Description of the 10 Indian cities (Delhi, Mumbai, Kolkata, Chennai, Bangalore, Hyderabad, Ahmedabad, Jaipur, Lucknow, Pune), their assigned crops, and the at-risk vs.\ control city grouping.

\subsection{Dataset 1: Daily Atmospheric Climate}
Placeholder: Open-Meteo / Kaggle Historical Weather Dataset covering 10 cities × 25 years of daily temperature, precipitation, and wind; humidity derived via Steadman apparent-temperature inversion.

\subsection{Dataset 2: CGWB Groundwater Levels}
Placeholder: Central Groundwater Board seasonal well readings (2000–2022) for 2,759 wells aggregated within 50 km city radii; imputation strategy for Jaipur (nearest-well proxy) and post-2022 gap.

\subsection{Dataset 3: FAO GAEZ v4 Crop Suitability}
Placeholder: GeoTIFF rasters at 9 km resolution providing rainfed baseline suitability scores (1–7) for each city-crop pair, extracted at city centroids.

% -----------------------------------------------------------------------
\section{Methodology}
\label{sec:methodology}
% -----------------------------------------------------------------------

\subsection{Köppen-Geiger Classification and Transition Detection}
Placeholder: Rule-based classification per Peel et al.\ (2007) applied annually; persistence filter requiring 3+ consecutive years before a transition is recorded.

\subsection{Annual Feature Engineering}
Placeholder: Derivation of 17 time-varying features (GDD accumulation, monsoon onset DOY, crop water deficit, recharge efficiency, dual-deficit flag, etc.) and 5 static features from ECOCROP and GAEZ.

\subsection{Crop Viability Loss Event (CVLE) Definition}
Placeholder: Formal dual-deficit definition requiring consecutive-year atmospheric and hydrological failure with at least 2 of 3 threshold breaches; pseudocode provided as Algorithm~\ref{alg:cvle}.

\begin{algorithm}
\caption{CVLE Label Construction}
\label{alg:cvle}
\begin{algorithmic}[1]
\FOR{each city}
  \FOR{$t = 1$ \TO $T$}
    \STATE $\textit{consec} \leftarrow \textit{dual\_deficit}[t] = 1 \;\wedge\; \textit{dual\_deficit}[t{-}1] = 1$
    \STATE $s \leftarrow \mathbb{1}[\textit{sowing\_window\_miss}[t] > 0.6] + \mathbb{1}[\textit{crop\_water\_deficit}[t] > 0.4] + \mathbb{1}[\textit{gdd\_adequate}[t] = 0]$
    \STATE $\textit{cvle\_label}[t] \leftarrow 1$ \textbf{if} $\textit{consec} \wedge s \geq 2$ \textbf{else} $0$
  \ENDFOR
\ENDFOR
\end{algorithmic}
\end{algorithm}

\subsection{Model Architectures}
Placeholder: Overview of all 8 models — Temporal Fusion Transformer (TFT), Random Forest, Logistic Regression, XGBoost, LightGBM, LSTM, TCN, and Vanilla Transformer — with hyperparameters.

\subsection{Train/Validation/Test Split}
Placeholder: Temporal split (train 2000–2018, val 2000–2021, test 2022–2024) with no data shuffling to prevent leakage; class imbalance (7.6\% positive) and handling strategies.

\subsection{Evaluation Metrics and Threshold Optimisation}
Placeholder: AUC-ROC, macro-F1, Brier score, accuracy; classification threshold selected by F1-macro grid search over the validation set rather than fixed at 0.5.

\subsection{Uncertainty Quantification}
Placeholder: Monte Carlo dropout (50 forward passes) for LSTM, TCN, and Transformer; tree variance for Random Forest; Expected Calibration Error (ECE) computed with 10-bin histogram.

\subsection{Feature Group Ablation}
Placeholder: Five mutually exclusive feature groups (Climate, Phenology, Hydrology, Static Context, All Features) evaluated independently on RF, XGB, and LSTM.

% -----------------------------------------------------------------------
\section{Experiments and Results}
\label{sec:results}
% -----------------------------------------------------------------------

\subsection{Unified Test-Set Performance}
Placeholder: Table comparing all 8 models on AUC-ROC, macro-F1, Brier score, accuracy, and optimal threshold (test set: $n=30$, years 2022–2024).

\subsection{Pairwise Statistical Significance}
Placeholder: Wilcoxon signed-rank test on per-sample Brier loss across 28 model pairs; results indicate which performance differences are statistically significant at $p < 0.05$.

\subsection{Feature Group Ablation Results}
Placeholder: Static context alone (GAEZ baseline + Köppen encoding) achieves AUC$\geq$0.98 for tree models; hydrological features are the dominant group for recurrent models.

\subsection{Uncertainty Quantification Results}
Placeholder: Random Forest tree-variance achieves the lowest ECE (0.045), outperforming all MC-dropout neural methods; interval widths and mean prediction std reported per model.

\subsection{SHAP Feature Importance}
Placeholder: Global mean $|\text{SHAP}|$ rankings from RF, XGB, and LGB; \texttt{gaez\_baseline\_class} and \texttt{koppen\_zone\_enc} dominate, confirming ablation findings.

\subsection{CVLE Event Timeline Per City}
Placeholder: Dated CVLE events for the 7 at-risk cities; control cities (Pune, Kolkata, Mumbai) produce near-zero CVLE rates, validating the dual-deficit formulation.

% -----------------------------------------------------------------------
\section{Discussion}
\label{sec:discussion}
% -----------------------------------------------------------------------

\subsection{Why Simple Models Outperform Deep Models}
Placeholder: The CVLE classification problem is largely linearly separable in feature space at $n=250$; LR achieves AUC$=$1.0 while TFT collapses to AUC$=$0.069.

\subsection{Why TFT Fails}
Placeholder: QuantileLoss is designed for multi-horizon regression, not binary classification; the 250-row dataset is insufficient for the ~115k-parameter model.

\subsection{Role of Static versus Dynamic Features}
Placeholder: Static GAEZ and Köppen features carry near-complete predictive signal; time-varying features add marginal value except for recurrent models exploiting groundwater trends.

\subsection{Groundwater as the Differentiating Signal for Recurrent Models}
Placeholder: Hydrological features achieve LSTM AUC$=$0.8621, the highest single-group AUC for sequential architectures, underscoring aquifer depletion as a lagged viability signal.

\subsection{Limitations and Validity Threats}
Placeholder: Key limitations include the small test set ($n=30$, 1 positive), CGWB data ending in 2022, Jaipur nearest-well proxy, humidity estimation from temperature inversion, and subjectively chosen CVLE thresholds.

% -----------------------------------------------------------------------
\section{Conclusion}
\label{sec:conclusion}
% -----------------------------------------------------------------------

VegShift introduces the first formal dual-deficit Crop Viability Loss Event framework tied to Köppen climate zone transitions, validated across 10 Indian cities over 25 years. Random Forest is the recommended deployment model (AUC$=$0.9655, ECE$=$0.045), while the CVLE definition is validated by near-zero false detection in stable control cities. Future work should extend to a wider geographic scope and incorporate ground-truth crop yield records for threshold calibration.

% -----------------------------------------------------------------------
\appendices
\section*{Acknowledgment}
Placeholder for acknowledgment of data sources (CGWB, Open-Meteo, FAO GAEZ) and funding.

% -----------------------------------------------------------------------
\begin{thebibliography}{00}
\bibitem{b1} Placeholder reference.
\end{thebibliography}

\EOD

\end{document}
